\subsection{11. Reproduktion und ehrliche
Liefergrenze}\label{reproduktion-und-ehrliche-liefergrenze}

\codeword{replay.py --frozen-only} führt die fünf Prüfer normal und mit
\codeword{-OO} aus, vergleicht beide JSON-Ergebnisse und die zwei
externen Originalberichte, prüft alle eingefrorenen Quellen und
protokolliert Änderungen der Originalorte gesondert. Ein Fehler führt zu
FAIL. Der zuvor aufgetretene Live-Quellenfehler bleibt in
\codeword{source\_drift\_failure\_receipt.json} sichtbar.

Die abgesicherten Zahlen stehen vollständig rational in
\codeword{two\_bank\_transfer\_normal.json}; Dezimalwerte dienen nur der
Lesbarkeit. Die analytische Beweiskette steht in Abschnitt 4. Das Paket
umfasst die unveränderten externen Arbeiten, deren Programme, die native
v1.6.4-Abhängigkeit und alle neuen Berichte. NumPy, SymPy und für das
ältere native Paket dessen separat dokumentierte Voraussetzungen sind zu
unterscheiden.

Die versionierte ausführliche Markdown-Lieferung enthält zusätzlich den
vollständigen eingefrorenen nativen Herleitungsbericht v1.6.4 als
historischen Anhang. Daneben gibt es ein kurzes Update und eine
bildliche Erklärung. \textbf{Die großen TFPT-Haupt-PDFs und die Webseite
werden in dieser Revision nicht als aktualisiert ausgegeben. Es gab
keinen Commit oder Push.}

\begin{center}\rule{0.5\linewidth}{0.5pt}\end{center}

\section{Nachprüfung der beiden zuletzt eingegangenen
Untersuchungen}\label{nachpruxfcfung-der-beiden-zuletzt-eingegangenen-untersuchungen}

Integrierter Nachtrag zu v1.6.5 · 15. September 2026

Quellen: die Anlagen \codeword{8e110ae1-19e2-4d47-bfc8-ebe4fe9ab81b} und
\codeword{b928b72a-6669-4c67-a75d-124095791c6f}, unverändert
eingefroren. Zusätzlich wurden die vier genannten Programme
\codeword{t1\_fixed.py}, \codeword{stabilizer.py},
\codeword{t5\_twobank.py}, \codeword{t5\_varresp.py} vollständig gelesen
und archiviert. Eine eigene Nachrechnung prüft die kritischen
Folgerungen; sie ist \textbf{kein vollständiger Replay aller sechs
Sonden} des fremden Arbeitsordners.

\subsection{1. Übersicht: übernehmen, korrigieren oder
offenlassen}\label{uxfcbersicht-uxfcbernehmen-korrigieren-oder-offenlassen}

{\def\LTcaptype{none} % do not increment counter
\begin{longtable}[]{@{}
  >{\raggedright\arraybackslash}p{(\linewidth - 2\tabcolsep) * \real{0.5000}}
  >{\raggedright\arraybackslash}p{(\linewidth - 2\tabcolsep) * \real{0.5000}}@{}}
\toprule\noalign{}
\begin{minipage}[b]{\linewidth}\raggedright
Neue Aussage
\end{minipage} & \begin{minipage}[b]{\linewidth}\raggedright
Ergebnis der Prüfung
\end{minipage} \\
\midrule\noalign{}
\endhead
\bottomrule\noalign{}
\endlastfoot
Innerer Viererzyklus erhält alle 60 W-Kanäle & Bestätigt, jetzt
einschließlich der richtigen Fermion-Vorzeichen und explizitem
Bosonlift \\
Grundzustandseindeutigkeit/Singulett/Gaplücke weiterhin unbekannt &
Überholter Stand: unter dem festgelegten nativen Vertrag bereits
bewiesen; voller Vektor weiter unbekannt \\
Leerer Grundzustand für alle μ≥0 & Falsch; bei μ=0 liegt der eindeutige
N=64-Grundzustand unter −1,129636Δ \\
Kommutant 7 als absolute Untergrenze & Nur im benannten blocktreuen
erweiterten Operationsvertrag; siehe Hauptbericht Abschnitt 8 \\
Beliebige Einteilchenmatrix ist damit als natives Fock-Wort verfügbar &
Nicht gezeigt; assoziative Matrixalgebra, Lie-Kontrolle und Fock-Lift
sind verschiedene Fragen \\
Kanal-Variationszustand hat E=−0,0196Δ & Die angegebene Phase im Code
ergibt gegen den gepinnten P=f\_jf\_i-Vertrag +0,05736Δ; ein relatives
Vorzeichen korrigiert es \\
99,892-\%-Transfer bei t≈2484 & Numerisches Maximum eines
Vierzustands-Paarmodells im untersuchten Zeitfenster; kein
erster/globaler exakter Maximalsatz und kein nativer Einloch-Transfer \\
Gleiche Zweigableitungen genau bei ε=0 & Algebraisch richtig; die
Eigenwertlücke bleibt dort √32· \\
SU(4)³-Anomalie =16 & Richtig unter A(4)=1 und der zusätzlichen
Interpretation aller (16,4) als gleichhändige Weylfelder; Eichung ist
eine weitere Voraussetzung \\
Gravitationsanomalie =64 & Keine reine perturbative Gravitationsanomalie
in 3+1D; 64 kann bei zusätzlich deklarierter gleichgeladener U(1) die
gemischte U(1)-Gravitationsanomalie zählen \\
Mit einer Ableitung kein (1,1)-Tensor bei Dimension ≤4 & Falsch: der
Energie-Impuls-Tensor ist ein Gegenbeispiel; daraus folgt aber noch kein
dynamisches Graviton \\
Kein lokaler kovarianter kinetischer Term für (1,0) & Zu pauschal: der
bilineare Ein-Ableitungs-Term fehlt, ein Zwei-Ableitungs-Skalar
existiert; Positivität/Constraints/Herkunft bleiben zu prüfen \\
\end{longtable}
}

„Negativ geschlossene Route`` bedeutet nicht „T2, T3, T4 oder T7
geschlossen``. Ein ausgeschlossener Ansatz beantwortet nicht die
jeweilige physikalische Existenzfrage. Die beiden gelieferten Texte
widersprechen sich vor allem beim bereits bewiesenen modellinternen
Grundzustand; sie dürfen nicht als gleichzeitig aktueller einheitlicher
Status zitiert werden.

\subsection{2. Positiver neuer Anschluss: der korrekt angehobene
Viererzyklus}\label{positiver-neuer-anschluss-der-korrekt-angehobene-viererzyklus}

Die Abbildung der Fermionmarken lautet \(p(4s+a)=4s+(a+1\bmod4)\). Auf
Paaren ist zwingend das Exteriorvorzeichen zu beachten: Falls
\(p(i)>p(j)\), erhält das sortierte Paar ein Minuszeichen. Der
gelieferte \codeword{stabilizer.py} lässt dieses Zeichen weg. Das wäre
im Allgemeinen falsch; für den untersuchten Viererzyklus überlebt das
positive Resultat jedoch auch den korrekten Test.

Alle 60 W-Zeilen werden mit ihren Vorzeichen wieder auf W-Zeilen
abgebildet. Wir haben den zugehörigen signierten Bosonoperator \(R_b\)
explizit gebildet und exakt geprüft:

\[\adjustbox{max width=\linewidth}{$\displaystyle 
 W'=R_bW,\qquad R_bR_b^T=I_{60},\qquad R_b^4=I_{60},
 \qquad R_b^2\ne I_{60}.
$}\]

Damit existiert ein \textbf{gemeinsamer nativer Tensor-Automorphismus}
auf Fermionen und Bosonen, nicht nur eine Übereinstimmung von
Zeilenzahlen. Die gleichzeitige Transformation erhält Paarwechselwirkung
und Bosonzahl. Sie beweist eine Ordnung-vier-Symmetrie, nicht deren
operative Verfügbarkeit.

Dieser Permutationszyklus ist nicht die zentrale Matrix \(iI_4\) von
SU(4). Seine Determinante auf dem Viererfaktor ist −1. Mit einer
zusätzlichen Ladungsphase kann man geeignete Gruppenlifts vergleichen;
die Permutation, die zentrale Phase, der Clock und die geometrische
Glue-Markierung sind dadurch aber nicht schon identifiziert. Genau diese
Intertwiner-Frage ist ein sinnvoller neuer Anschluss, statt noch einmal
nur vier zu zählen.

\subsection{3. Warum der Einteilchen-Abschluss den Fock-Operationssatz
nicht
schließt}\label{warum-der-einteilchen-abschluss-den-fock-operationssatz-nicht-schlieuxdft}

Die Irreduzibilität der Einteilchendarstellung kann ihre assoziative
Matrixalgebra zu \(B(\mathbb C^{64})\) machen. Daraus folgt nicht, dass
jede Matrix als ein zulässiges physisches Wort verfügbar ist, und auch
nicht, dass ihre zweite Quantisierung bereits erzeugt wird.

Ein kleinstes exaktes Gegenbeispiel zur falschen Liftregel:
\(A=|1\rangle\langle1|\), \(B=|2\rangle\langle2|\) auf zwei Moden. Dann
\(AB=0\), also \(d\Gamma(AB)=0\), aber
\(d\Gamma(A)d\Gamma(B)=n_1n_2\ne0\). Der zweite Quantisierungsschritt
ist ein Lie-, nicht ein assoziativer Algebra-Homomorphismus dieser Art.

Auch \(\operatorname{diag}(i,1,1,1)\) und
\(\operatorname{diag}(i,i,1,1)\) haben Determinanten i beziehungsweise
−1. Ihre Zugehörigkeit zur \textbf{linearen Matrixalgebra} aus
SU(4)-Generatoren und Identität ist kein exaktes SU(4)-Gruppenwort.
Projektive Gleichheit oder eine zusätzliche U(1)-Phase kann helfen, muss
aber mit der Bosonwirkung, Ladung und Referenz gemeinsam ausgewiesen
werden.

Der fremde Code folgert außerdem die Irreduzibilität des
37.888-dimensionalen dunklen Dreifermionraums aus der vollen
Einteilchenmatrixalgebra. Diese Schlussregel ist nicht begründet.
Bereits die im anderen neuen Text angegebenen drei dunklen irreduziblen
Typen widersprechen der pauschalen Eins-Block-Lesart unter der bloßen
Quellsymmetrie. Die Zahl 8.732.673 mag als sehr schwache obere Schranke
anderweitig verträglich sein; die im Code angegebene Herleitung belegt
sie nicht. Die SVD-/Toleranzprüfungen des Codes sind zudem numerisch,
nicht allein wegen „PASS`` exakte Lie-Beweise.

\subsection{4. Zustandswahl: keine Rückkehr vor den abgesicherten
Grundsatz}\label{zustandswahl-keine-ruxfcckkehr-vor-den-abgesicherten-grundsatz}

Der vorhandene native Satz beweist Eindeutigkeit, N=64, Singulett und
positive Lücke bei μ=0 im angegebenen Kopplungsbereich. Eine komplette
Clebsch-Gordan-Serie oder Voll-Diagonalisierung ist dafür nicht nötig;
Vergleichsungleichungen und Minmax waren gerade der einfache Ausweg. Die
physische Auswahl von H beziehungsweise μ=0 bleibt eine andere Frage.

Das grobe Sandwich \([-1.2,-.0196]\Delta\) ist nach einer
Phasenkorrektur zwar verträglich, aber wesentlich schwächer als
\((-1.158089,-1.129636)\Delta\). Es ist keine Verschärfung.

Im tatsächlich angegebenen Variationscode werden die Fermionen in der
Reihenfolge j, dann i gelöscht. Das liefert \(f_if_jF=-f_jf_iF\),
während der native Vertrag \(P_A=\sum W_{A,ij}f_jf_i\) verwendet. Bei
unverändertem \(g=+\Delta/20\) und den angegebenen
Variationskoeffizienten folgt deshalb

\[\adjustbox{max width=\linewidth}{$\displaystyle 
 E_{\rm Code}/\Delta=\frac12-\frac{23\sqrt3}{90}
 =+.057364793621\ldots,
$}\]

nicht der behauptete negative Wert. Mit korrigierter relativer Phase
lautet er \(1/2-3\sqrt3/10=-.019615242271\ldots\). Alternativ wäre eine
konsequente andere P/g-Phasenkonvention möglich; die Quelle muss sie
dann überall führen. Die gleichzeitige Einteilchendichte des einfachen
Variationszustands bleibt von dieser Phasenreparatur unberührt, weil
seine Bosonzahlkomponenten orthogonal sind. Sie ist keine dynamische
Greenfunktion auf dem nativen Grundzustand und ersetzt dessen schon
kontrollierte Antwort nicht.

\subsection{5. Feldtheorie: falsche Verbote entfernen, echte Bedingungen
behalten}\label{feldtheorie-falsche-verbote-entfernen-echte-bedingungen-behalten}

\subsubsection{5.1 Anomalien}\label{anomalien}

Für die zusätzlich angenommene 3+1D-Weylinterpretation lautet das lokale
Anomaliepolynom bis auf Konventionsvorzeichen

\[\adjustbox{max width=\linewidth}{$\displaystyle 
 I_6=[\widehat A(T)\,\mathrm{ch}_R(F)]_6
 =\mathrm{ch}_3(F)-\frac{p_1(T)}{24}\,\mathrm{ch}_1(F).
$}\]

Der reine gravitative Anteil \([\widehat A]_6\) verschwindet; seine
Formgrade sind Vielfache von vier. Das ist die bekannte dimensionale
Unterscheidung bei
\href{https://collaborate.princeton.edu/en/publications/gravitational-anomalies/}{Álvarez-Gaumé
und Witten}. Eine Spur-/Weylanomalie ist nochmals ein anderer Begriff.

Auf (16,4) ist der SU(4)-Kubikkoeffizient 16 bei Normierung A(4)=1. Der
exakte Test \(t=\operatorname{diag}(1,1,1,-3)\) hat
\(\operatorname{tr}t=0\), \(\operatorname{tr}t^3=-24\). Ist SU(4)
\textbf{dynamisch geeicht}, muss die vollständige Theorie diese
Eichanomalie kompensieren. Als globale Symmetrie kann sie eine
't-Hooft-Anomalie tragen; dann folgt keine pauschale Spiegelpflicht. Ein
vollständiger konjugierter Spiegel ist ein möglicher Ausgleich, aber
hier nicht als einzige oder allgemein minimale Lösung bewiesen.

„Gravitativ 64`` kann sinnvoll eine \textbf{gemischte
U(1)-Gravitationsanomalie} meinen, wenn alle 64 linksgetragenen
Weylfelder explizit U(1)-Ladung eins erhalten. Das muss samt
Eich-/Globalstatus gesagt werden. Für SU(4) allein ist der gemischte
lineare Spurkoeffizient null. Die Quellenmarke N darf nicht ohne Beweis
in eine dynamisch geeichte chirale U(1) umgedeutet werden.

\subsubsection{5.2 Der behauptete Spin-2-Ausschluss ist
falsch}\label{der-behauptete-spin-2-ausschluss-ist-falsch}

Die Lorentzdarstellungen liefern exakt

\[\adjustbox{max width=\linewidth}{$\displaystyle 
 (\tfrac12,\tfrac12)\otimes(\tfrac12,\tfrac12)
 =(0,0)\oplus(1,0)\oplus(0,1)\oplus(1,1).
$}\]

Eine Ableitung des Vektorbilinears \(\psi^\dagger\bar\sigma_\mu\psi\)
kann daher einen symmetrischen spurfreien (1,1)-Tensor bilden.
Insbesondere hat der Energie-Impuls-Tensor
\(T_{\mu\nu}\sim i\psi^\dagger\bar\sigma_{(\mu}
\overleftrightarrow\partial_{\nu)}\psi\), mit passender Spurbehandlung,
Dimension \(3/2+3/2+1=4\). Ein Energie-Impuls-Tensor kann bereits in
einer festen Hintergrundraumzeit definiert werden; dynamische
Diffeomorphismen müssen nicht vorausgesetzt werden, um dieses
Gegenbeispiel zu formulieren.

Das schließt \textbf{T7 nicht}. Ein vorhandener Tensor ist kein
masseloser Spin-2-Pol. Das sinnvolle Ziel ist später sein transversaler,
spurfreier Korrelator auf demselben räumlichen Parent, mit positiver
Norm, beiden Helizitäten und universeller Kopplung. Wir entfernen hier
eine falsche Darstellungsobstruktion, nicht die dynamischen
Nachweispflichten.

\subsubsection{5.3 (1,0)-Kinetik ist nicht generell
verboten}\label{kinetik-ist-nicht-generell-verboten}

Der bilineare Ein-Ableitungs-Term mit einem (1,0)-Feld und seinem
Adjungierten besitzt keinen Lorentzskalar. Mit zwei Ableitungen
existiert aber etwa

\[\adjustbox{max width=\linewidth}{$\displaystyle 
 \mathcal L_2\propto
 (\partial^{\alpha\dot\alpha}\Phi_{\alpha\beta})
 (\partial^{\beta\dot\beta}\bar\Phi_{\dot\alpha\dot\beta})
$}\]

für symmetrisches \(\Phi\). Alle Spinorindizes sind kontrahiert. Bei
rein zeitartigem Impuls ist sein Symbol proportional zu
\(\omega^2(|\Phi_{11}|^2+2|\Phi_{12}|^2+|\Phi_{22}|^2)\), also nicht
identisch null. Das ist ein Gegenbeispiel zum uneingeschränkten Satz
„kein lokaler kovarianter kinetischer Term``. Es beweist \textbf{nicht}
Positivität des vollen Hamiltonoperators, korrekte Constraints,
gewünschte Freiheitsgrade oder eine native Quelle. Diese bleiben die
entscheidenden Tests.

\subsection{6. Die Vierzustandsrechnung richtig
lesen}\label{die-vierzustandsrechnung-richtig-lesen}

Das Modell mit \(s=\sqrt8g\), Bosonverbindung η und Diagonalen (0,Δ,Δ,0)
erlaubt kohärenten \textbf{Paartransport}. Die exakte Identität
\((H^3)_{41}=8g^2\eta\) stimmt. Unsere unabhängige Fließkomma-Abtastung
reproduziert auf \([0,3000]\) mit Schrittweite .05 den größten
gefundenen Wert .998920021869 bei 2484.05. Der erste
Gitter-Lokalhöchstwert liegt bereits bei 6.05. Weder „erster
Maximalzeitpunkt`` noch ein globales exaktes Maximum ist damit bewiesen.
Diese kleine Rechnung wird nicht mit dem vollständigen Einloch-Satz des
Hauptberichts verwechselt.

Für \(F_\pm(\epsilon)=(\epsilon\pm\sqrt{\epsilon^2+32g^2})/2\) gilt
\(F'_+-F'_-=\epsilon/\sqrt{\epsilon^2+32g^2}\). Bei ε=0 stimmen die
Ableitungen überein, aber die \textbf{Eigenwertlücke} beträgt
\(\sqrt{32}|g|>0\). Ein verschwindender nackter Vermittlerparameter und
eine verschwindende wechselwirkende Anregungslücke sind verschieden. Ein
allgemeiner relativistischer Kegelsatz oder ein zwingendes neues
Renormierungsprogramm folgt aus dieser Ableitungsidentität nicht.

\subsection{7. Konsequenz für die nächste
Forschungsrevision}\label{konsequenz-fuxfcr-die-nuxe4chste-forschungsrevision}

Die sinnvollen positiven Anschlüsse sind jetzt schärfer:

\begin{enumerate}
\def\labelenumi{\arabic{enumi}.}
\tightlist
\item
  Den expliziten Ordnung-vier-Lift mit dem tatsächlich markierten Clock
  und der Herkunft des A3-Index vergleichen: vollständige Intertwiner,
  Ladungsphase und Bosonwirkung, nicht nur Gruppennamen.
\item
  Einen Quellenoperator finden, der die lokale Parität beider
  operational bestimmten Teile gleichzeitig ändert. Die native
  Z4-Symmetrie selbst tut das nicht. Der kontrollierte Einloch-Transfer
  ist bereits als Akzeptanztest für einen solchen Operator verfügbar.
\item
  Den Zwei-Ableitungs-Feldkandidaten auf Positivität und Constraints
  prüfen, bevor er als relativistischer Adapter zählt. Ein bloßes
  Nichtnull-Symbol ist noch keine gesunde Feldtheorie.
\item
  Den Energie-Impuls-Tensor als zulässigen \textbf{Kandidaten} behalten,
  aber erst auf dem gemeinsam konstruierten räumlichen Träger nach einem
  dynamischen Spin-2-Pol suchen. Kein vorhandener Spin-2-No-go
  rechtfertigt hier das Aufgeben dieses Anschlusses.
\end{enumerate}

62 zusätzliche Prüfbedingungen bestehen normal und optimiert identisch.
Exakte Identitäten, bedingte Darstellungsargumente und die numerische
Zeitfenstersuche sind im Ergebnisbericht getrennt. Der externe Status
„alle Restfragen geschlossen oder auf genau ein Stück reduziert`` wird
nicht übernommen. Die vollständige TOE bleibt offen.

\begin{center}\rule{0.5\linewidth}{0.5pt}\end{center}

\section{Historischer Beweisanhang: vollständiger nativer
Herleitungsstand
v1.6.4}\label{historischer-beweisanhang-vollstuxe4ndiger-nativer-herleitungsstand-v1.6.4}

Dieser Anhang bewahrt die frühere Herleitung vollständig. Seine
damaligen Statussätze werden durch die aktuelle Konsolidierung und den
Nachtrag oben ergänzt; er ist kein ungeprüft übernommener neuer
Gesamtstatus. Die externe laufende Symmetrieergänzung wurde separat
geprüft und ist nicht unbemerkt in diesen eingefrorenen Anhang
eingegangen.

\section{TFPT / Universalraum: nativer Pol, Bewegungsgleichung und
minimale
Feldtypen}\label{tfpt-universalraum-nativer-pol-bewegungsgleichung-und-minimale-feldtypen}

\textbf{Konsolidierte Forschungsfortsetzung v1.6.4 · 15. September 2026}

Diese Revision verbindet den während der Arbeit neu eingegangenen
Polsatz v1.6.3 mit einer unabhängig begonnenen Untersuchung der nativen
Bewegungsgleichung. Sie ergänzt die bisherigen Hauptdokumente; sie ist
keine verkürzte Neufassung des vollständigen Hauptbuchs. Die vorhandenen
Haupt- und Update-PDFs wurden in dieser Runde nicht verändert.

